\par\Needspace{6\baselineskip}\medskip
\hypertarget{native-126-article093}{}\pdfbookmark[2]{93.}{article-093}
\begin{center}93.\end{center}
Um die Lösung der Gleichung 15 in den praktisch häufigsten Fällen möglichst bequem zu machen, fügen wir am Ende dieses Werkes eine besondere Tafel (Tafel II) bei. Sie gibt für Werte von $h$ von 0 bis 0,6 die entsprechenden Logarithmen von $yy$ an, die mit größter Sorgfalt auf sieben Dezimalstellen berechnet sind. Das Argument $h$ wächst von 0 bis 0,04 in Schritten von einem Zehntausendstel, wodurch die zweiten Differenzen von $\log yy$ zum \SourceLeaf{127}Verschwinden gebracht werden. In diesem Teil der Tafel genügt daher die einfache Interpolation. Da die Tafel bei gleicher Schrittweite über den gesamten Bereich sehr umfangreich geworden wäre, musste sie von $h=0{,}04$ bis zum Ende in Schritten von nur einem Tausendstel fortgeführt werden. In diesem letzten Teil muss man daher die zweiten Differenzen berücksichtigen, wenn man Fehler von einigen Einheiten in der siebten Stelle vermeiden will. In der Praxis kommen übrigens kleinere Werte von $h$ bei Weitem am häufigsten vor.

Sowohl Gleichung 15 bei Überschreitung der Tafelgrenze durch $h$ als auch Gleichung $15^{\ast}$ lassen sich ohne Schwierigkeit nach der indirekten Methode oder nach anderen hinreichend bekannten Methoden lösen. Es sei noch bemerkt, dass ein kleiner Wert von $g$ nur bei sehr exzentrischen Bahnen mit einem negativen Wert von $\cos f$ vereinbar ist, wie sich unmittelbar aus Gleichung 20 ergibt, die weiter unten in Art. 95 angegeben wird\footnote{Aus dieser Gleichung geht hervor, dass für negatives $\cos f$ der Winkel $\varphi$ jedenfalls größer sein muss als $90^0-g$.}.
